\chapter{Introduction}
\label{chap:introduction}

\section{Adversarial Robustness and the Cross-Attack Gap}

Modern neural networks achieve high accuracy on standard image classification
benchmarks, yet they can be sensitive to small, carefully crafted perturbations
of their inputs.  These adversarial perturbations are often imperceptible to
human observers, but they can change the model's prediction with high
reliability.  The phenomenon has motivated a large body of work on adversarial
training, which augments the training process with adversarial examples so that
the resulting classifier learns to resist such perturbations.

Adversarial training has become the dominant empirical defense against
adversarial attacks.  However, the effectiveness of adversarial training
typically depends on the attack used during training.  Different attacks employ
different perturbation norms, optimization objectives, and step-size schedules,
and a model trained against one attack may not transfer its robustness to
attacks with different characteristics.  For example, a model that is robust
under a projected gradient descent attack in the $\ell_\infty$ norm may remain
vulnerable to optimization-based attacks in the $\ell_2$ norm, and vice versa.
This creates what this report calls the \emph{cross-attack robustness gap}: the
difference in a model's robustness when evaluated under attacks that differ from
the one used during training.

The cross-attack gap is not merely a theoretical concern.  Evaluation suites
that include attacks from multiple norms, step counts, and optimization
objectives routinely show that single-attack adversarial training produces
models whose robustness is concentrated on the source threat family.  Once the
evaluation moves outside that family, performance can degrade substantially.
This observation suggests that robustness should be measured not by a single
attack score but by how well a model performs across a diverse set of threat
models.

One natural response is to train the model against multiple attacks
simultaneously.  Multi-attack adversarial training constructs adversarial
examples from more than one source attack during each training step and averages
the resulting losses.  This strategy provides broader coverage than single-attack
training, but it treats all attack-induced examples uniformly.  If some attacks
or attack-induced subgroups are systematically harder than others, uniform
averaging can hide those failure modes behind an improved overall mean.

This report focuses on improving the balance of robustness across a fixed attack
suite rather than claiming universal robustness against arbitrary perturbations.
The experimental scope is bounded to CIFAR-10 image classification with a
defined set of evaluation attacks, and the claims are conditioned on the
specific training and evaluation protocol described in later chapters.  This
motivates the central question of this report: whether adversarial training can
be made less dependent on a single source attack by treating attacks and
attack-induced failures as structured domains.

\section{Stage 1 Findings and Their Limitations}

The present work continues from a Stage~1 investigation that diagnosed
attack-dependent robustness under single-attack adversarial training.  Stage~1
trained classifiers using individual source attacks and evaluated them against a
broader attack suite.  The main observation was that single-attack training
tends to specialize: a model trained against an $\ell_\infty$ source attack
achieved strong robustness under held-out $\ell_\infty$ attacks but weaker
performance under $\ell_2$ attacks, and symmetrically for an $\ell_2$ source
attack.

Stage~1 also established an evaluation discipline that carries forward into the
present report.  The attack suite, the use of repeated random seeds, and the
separation of source and held-out attacks were all informed by Stage~1
observations.  However, Stage~1 was primarily evaluative and diagnostic: it
measured the cross-attack gap but did not propose a training method to reduce
it.

The limitation of Stage~1 is therefore clear.  Diagnosis alone does not produce
a training intervention.  Knowing that single-attack robustness is
attack-dependent motivates the search for a training strategy that balances
robustness across attacks, but the strategy itself must be designed, implemented,
and tested.  Stage~2 addresses this gap by introducing a progression of
multi-attack and group-aware training methods, evaluated under a controlled
experimental protocol.

The present report therefore moves from diagnosis to intervention.  It takes the
cross-attack gap observed in Stage~1 as its starting point and develops a
method that treats attack-induced adversarial examples as domain-like groups.
The goal is to shift training emphasis toward harder groups without abandoning
the stability of a uniform multi-attack baseline.

\section{Problem Statement}

The problem studied in this report is supervised image classification on
CIFAR-10, where a classifier is trained under adversarial examples generated by
two source attacks: PGD-$\ell_\infty$ and DDN-$\ell_2$.  Each source attack
induces a different distribution of adversarial examples, and these distributions
can be viewed as domain-like groups with distinct loss landscapes and failure
patterns.  The trained model is then evaluated under a broader suite of eight
attacks spanning both the $\ell_\infty$ and $\ell_2$ families, along with a
separate AutoAttack-$\ell_\infty$ stress test.

Uniformly averaging the losses across source attacks during training can hide
hard attacks, hard classes, or latent subgroups that consistently receive higher
loss.  If certain attack-induced regions are systematically underweighted, the
resulting model may appear robust on average while remaining vulnerable in
specific evaluation cells.  This motivates the need for a training objective
that can detect and emphasize difficult groups during optimization.

The problem studied in this report is not to solve adversarial robustness in
general, but to improve the balance of robustness across a predefined attack
suite under a reproducible training and evaluation protocol.  The goal is to
improve average robustness across seen and held-out attacks while avoiding severe
degradation on the hardest evaluation cases, all within the computational and
dataset constraints of CIFAR-10 and ResNet-18.

\section{Research Questions and Hypotheses}

This report is organized around four research questions, each targeting a
different aspect of cross-attack robustness.  These questions are designed so
that the experimental evidence in Chapter~\ref{chap:results_analysis} can
address them directly, and Chapter~\ref{chap:conclusion} can synthesize the
answers.

\textbf{RQ1: How large is the cross-attack robustness gap under single-attack
training?}  The hypothesis is that single-attack adversarial training will
produce attack-dependent robustness profiles, with each single-attack method
protecting its own source family more strongly than attacks outside that
training objective.

\textbf{RQ2: Does multi-attack training improve robustness across seen and
held-out attacks?}  The hypothesis is that Multi-AT, a uniform multi-attack
adversarial training baseline, will improve robustness balance compared with
single-attack baselines by covering both source threat geometries, though
uniform averaging may not fully remove hard attack-specific or class-specific
failure modes.

\textbf{RQ3: Can group-aware training improve robustness and stability beyond
uniform multi-attack training?}  The hypothesis is that AttackDRO++, which
discovers latent hard groups and upweights them during training, will improve
average robustness and training stability over Multi-AT in the main setting.
However, hardest-case robustness and transfer behavior should be interpreted
conservatively, as the advantage may be more limited under those evaluation
conditions.

\textbf{RQ4: Which design choices are responsible for the final behavior?}  The
hypothesis is that moderate cluster count, moderate anchor strength, and a
stable cluster-refresh schedule should produce better robustness balance than
extreme settings, providing practical configuration guidance for the method.

\section{Contributions of This Work}

This report makes the following contributions within the bounded experimental
setting of CIFAR-10 and ResNet-18:

\begin{enumerate}
\item It proposes an \emph{attacks-as-domains} framing for cross-attack
robustness, in which adversarial examples generated by different source attacks
are treated as domain-like groups whose losses can be balanced during training.

\item It implements a progression of training methods: a uniform Multi-AT
baseline, AttackDRO with group distributionally robust optimization (DRO) over
source attack identities, and AttackDRO++ with group DRO over discovered latent
clusters.

\item It introduces gradient fingerprint clustering features that augment the
standard penultimate representation with projected classifier-head gradients,
enabling the grouping mechanism to reflect optimization-level difficulty rather
than representation similarity alone.

\item It introduces a uniform-anchored training objective that keeps the
AttackDRO++ objective tied to the Multi-AT training signal while still allowing
difficult discovered groups to receive additional emphasis.

\item It evaluates the proposed method through a reproducible protocol that
includes aggregate whitebox comparisons, per-attack and per-class drilldowns,
graybox transfer analysis across independently trained surrogate models,
statistical significance testing with paired seeds, and hyperparameter
sensitivity studies.
\end{enumerate}

\section{Scope and Limitations}

The experimental scope of this report is bounded in several ways that should be
stated before the method and results chapters.

The dataset used throughout is CIFAR-10, a standard benchmark for adversarial
training research but a small, low-resolution image dataset.  The main
architecture is ResNet-18, adapted for CIFAR-10 resolution.  All central claims
are conditioned on this dataset and architecture combination.

The source attacks used during training are PGD-$\ell_\infty$ and
DDN-$\ell_2$.  The evaluation suite includes eight attacks spanning both
$\ell_\infty$ and $\ell_2$ families, plus a separate
AutoAttack-$\ell_\infty$ stress test that is reported independently from the
main aggregate metrics.  The evaluation does not include certified robustness,
full adaptive robustness, or robustness against attacks outside the defined
suite.

Graybox transfer evaluation is included, using independently trained surrogate
models from the same method families.  This tests whether robustness behavior
remains stable under transferred attacks, but it is still bounded to the chosen
model families and attack suite.  Supplementary architecture evidence from a
single WRN-28-10 run is provided in Appendix~\ref{app:extras}, but it is not
part of the main paired claims and does not support a broad architecture-level
generalization.

This report does not claim universal robustness.  It does not establish that the
proposed method generalizes to other datasets, architectures, or attack
configurations beyond those tested.  The conclusions are intended to be useful
within the defined protocol and to motivate further investigation under broader
conditions.

\section{Report Organization}

The remainder of this report is organized as follows.
Chapter~\ref{chap:background} provides the background needed to understand the
proposed method, covering adversarial attacks, adversarial training, distributionally
robust optimization (DRO), domain generalization, clustering,
and the statistical tools used for multi-seed evaluation.
Chapter~\ref{chap:related_work} reviews related work on multi-attack adversarial
training, group-aware robustness, domain generalization, and cluster-based
training, and identifies the remaining gap that this report addresses.

Chapter~\ref{chap:methodology} presents the proposed methodology in a
progression from the attacks-as-domains formulation through Multi-AT, AttackDRO,
and AttackDRO++, including gradient fingerprint clustering and the
uniform-anchored objective.  Chapter~\ref{chap:experimental_setup} defines the
experimental protocol: dataset, architecture, attack suite, training
configuration, hyperparameters, statistical significance protocol, evaluation
metrics, and graybox transfer design.

Chapter~\ref{chap:results_analysis} presents the results and analysis, moving
from aggregate whitebox robustness to class-wise drilldowns, graybox transfer,
and hyperparameter sensitivity studies.  Chapter~\ref{chap:conclusion} concludes
the report with answers to the research questions, a discussion of limitations,
and directions for future work.  The report therefore moves from motivation and
background to method design, then to a controlled empirical evaluation and a
final discussion of what the evidence supports.
